\documentclass{article}
\usepackage[margin=1in]{geometry}
\usepackage{amsmath,amssymb}
\usepackage[most]{tcolorbox}
\usepackage[hidelinks]{hyperref}

\newcommand{\BookName}{Pattern Recognition and Machine Learning}
\newcommand{\BookAuthor}{Christopher M. Bishop}
\newcommand{\ChapterName}{Introduction}
\newcommand{\ExerciseID}{1.6}

\title{Chapter: \ChapterName}
\author{Ahonakpon Guy Baguidi}
\date{}

\begin{document}

\maketitle

\begin{center}
  \large\textbf{Exercise \ExerciseID}

  \vspace{0.5em}

  \normalsize\BookName\ - \BookAuthor
\end{center}

\begin{tcolorbox}[breakable,colback=gray!5,colframe=gray!60,title=Solution]
In this exercise, we have to show that if two variables $x$ and $y$ are independent,
then their covariance is zero.
\par\medskip
Let $x$ and $y$ be two random variables defined on a probability space with joint probability 
density function $p(x, y)$. The covariance of $x$ and $y$ is defined as:
\begin{align*}
  \text{Cov}[x, y] &= \mathbb{E}[(x - \mathbb{E}[x])(y - \mathbb{E}[y])] \\
  &= \mathbb{E}_{xy}[xy] - \mathbb{E}[x]\mathbb{E}[y]
\end{align*}

If $x$ and $y$ are independent, then their joint probability density function can be 
written as the product of their marginal probability density functions:
\begin{align*}
  p(x, y) = p(x) p(y)
\end{align*}

Now we compute the expectation of the product $xy$. The form of the expectation depends
on the type of variables:
\begin{itemize}
  \item for discrete variables, expectation is written as a sum;
  \item for continuous variables, expectation is written as an integral.
\end{itemize}
In both cases, the same idea is used: because $x$ and $y$ are independent, the joint
distribution factors as $p(x,y)=p(x)p(y)$, and then we can separate the double sum or
double integral by using 	\textit{Fubini's theorem} (or Tonelli's theorem when the functions
are nonnegative).
\begin{align*}
  \mathbb{E}_{xy}[xy] &= \int \int xy p(x, y) dx dy \\
  &= \int \int xy p(x) p(y) dx dy \\ \quad\text{(by independence)} \\
  &= \left( \int x p(x) dx \right) \left( \int y p(y) dy \right) \\
  &= \mathbb{E}[x] \mathbb{E}[y]
\end{align*}

So the same proof works for discrete sums and continuous integrals; only the notation changes.
The important point is that the joint probability factors, and then the sum or integral can
be separated into one part depending only on $x$ and one part depending only on $y$.
\par\medskip

By substituting this result back into the expression for covariance, we have:

\begin{align*}
  \text{Cov}[x, y] &= \mathbb{E}_{xy}[xy] - \mathbb{E}[x]\mathbb{E}[y] \\
  &= \mathbb{E}[x] \mathbb{E}[y] - \mathbb{E}[x]\mathbb{E}[y] \\
  &= 0
\end{align*}

This shows that if two variables $x$ and $y$ are independent, then their covariance 
is zero, which means that there is no linear relationship between them. 
However, it is important to note that the converse is not necessarily true: 
if the covariance of two variables is zero, it does not imply that they are independent, 
as there could still be a nonlinear relationship between them.

\paragraph{Personal note on Fubini's theorem:} The key step in this derivation is not integration by parts. 
It is the use of Fubini's theorem, which lets us rewrite a double integral as an iterated integral when the function is integrable. 
For discrete random variables, the same idea is written with double sums instead of double integrals. 
Because independence gives $p(x,y)=p(x)p(y)$, the expression becomes a product of an $x$-part and a $y$-part:
\begin{align*}
\int \int xy \, p(x) p(y) \, dx dy = \left( \int x p(x) dx \right) \left( \int y p(y) dy \right)
\end{align*}
So this is really a theorem about changing the order of summation or integration, not about differentiating products.

\paragraph{Resources to learn more about Fubini's theorem:}
\begin{itemize}
  \item Rudin, Walter. \textit{Real and Complex Analysis}. McGraw-Hill, 3rd ed. 1987. Chapter 8, Theorem 8.8.
  \item \url{https://en.wikipedia.org/wiki/Fubini\%27s_theorem}.
\end{itemize}

\end{tcolorbox}

\end{document}
